% P5 iter-105 — live manifest per-field per-value coverage audit
\subsection{Exhibit: Per-field per-value coverage at $n{=}98$ manifests}
\label{sec:p5-iter105-coverage}

To stress-test the MIN-REPORT standard at the level of \emph{individual
declared-field values}, we re-audited every manifest under
\texttt{platform_hybrid/experiments/results/mega\_20260704/manifests/} ($n{=}98$ cells)
and classified each of the eight top-level manifest keys per cell into
one of five value-categories: \texttt{PRESENT\_PATH} (string starting with
\texttt{/}), \texttt{PRESENT\_KL\_CONCRETE} (string starting with
\texttt{kl-}), \texttt{PRESENT\_NA} (literal \texttt{n/a} or
\texttt{n/a-}* sentinel), \texttt{PRESENT\_KEYWORD} (any other concrete
string), \texttt{PRESENT\_NULL} (JSON \texttt{null}), \texttt{MISSING}
(key absent from the manifest). Analysis script:
\texttt{platform_modal/scripts/p5p8/p5\_iter105\_live\_field\_coverage.py}; bootstrap
$B{=}2000$, seed $20260705$. Per-cell classification:
\texttt{platform_hybrid/experiments/results/p5p8/p5\_iter105\_per\_field\_class.tsv}
($784{=}98{\times}8$ rows); aggregate:
\texttt{p5\_iter105\_per\_field\_summary.tsv}; unique-value inventory:
\texttt{p5\_iter105\_unique\_values.tsv} ($209$ rows); Item-2
\texttt{ref\_policy\_kl} cell-level inventory:
\texttt{p5\_iter105\_item2\_kl\_inventory.tsv}; machine summary:
\texttt{p5\_iter105\_summary.json}.

\begin{table}[t]
\centering
\caption{Per-value-category coverage of the eight declared manifest keys
across the $n{=}98$ live mega-campaign manifests. All eight keys are
present in $98/98$ cells (no \texttt{MISSING}); but the value-category
distribution splits cleanly into \emph{discriminative} (cell_id,
per\_step\_zvf\_path, group\_size\_schedule, heldout\_split,
decontamination\_notes) and \emph{constant-or-sentinel} (loss\_form,
ref\_policy\_kl, sampler\_backend\_precision) buckets.}
\label{tab:p5-iter105-coverage}
\small
\begin{tabular}{@{}lrrrrr@{}}
\toprule
field & path & KL & n/a & keyword & uniq \\
\midrule
cell\_id                     &   0 &   0 &  0 &  98 &  98 \\
loss\_form                   &   0 &   0 & 98 &   0 &   1 \\
ref\_policy\_kl               &   0 &   0 & 98 &   0 &   1 \\
sampler\_backend\_precision   &   0 &   0 &  0 &  98 &   1 \\
per\_step\_zvf\_path          &  98 &   0 &  0 &   0 &  98 \\
group\_size\_schedule         &   0 &   0 &  0 &  98 &   5 \\
heldout\_split               &   0 &   0 &  0 &  98 &   3 \\
decontamination\_notes       &   0 &   0 &  0 &  98 &   2 \\
\bottomrule
\end{tabular}
\end{table}

\textbf{Falsifiable headline H1 --- \texttt{PRESENT\_KL\_CONCRETE}
fraction on \texttt{ref\_policy\_kl} is exactly $0/98$ (CI $[0.000,0.000]$).}
No live corpus cell provides a concrete \texttt{kl-}* value; the field is
uniformly \texttt{n/a} because the corpus is sampling-only with no KL
regulariser. The Item-2 validation gap from Exhibit~6 is therefore
\emph{not} an emission bug --- the literal-sentinel string is a valid
\emph{declaration-of-absence}, but it is \emph{indistinguishable} from an
emission bug under naive parsing. Iter-105 makes the indictment sharp:
$98$ of $98$ cells carry a value, but $0$ of $98$ cells carry an
information-bearing value on Item 2.

\textbf{Falsifiable headline H2 --- $5/8$ declared fields are
stack-discriminative ($n_{\text{unique}} \ge 2$), $3/8$ are
stack-constant.}
The five discriminative fields (cell\_id, per\_step\_zvf\_path,
group\_size\_schedule, heldout\_split, decontamination\_notes) encode
$2 \cdot 3 \cdot 5 \cdot 2 \cdot 2 = 120$ axis-combinations, of which
$98$ are realised. The three constant-or-sentinel fields
(loss\_form, ref\_policy\_kl, sampler\_backend\_precision) encode no
stack information; they are \emph{audit primitives} (e.g.\ confirming
the corpus is sampling-only) rather than \emph{stack axes}. The split
sharpens the iter-97 schema audit: \emph{key presence} ($8/8$ per
cell) overstates \emph{stack-discriminative coverage} ($5/8$ per cell).
The naive ``$100\%$ present'' headline of Exhibit~6 is correct only if
constant sentinels count as items.

\textbf{Falsifiable headline H3 --- per-axis uniqueness profile
matches \texttt{cells.tsv} $5$-axis factorial structure.}
unique-value counts on the discriminative fields ($98, 98, 5, 3, 2$)
factor as $98 = 1 \cdot 1 \cdot 5 \cdot 2 \cdot 2 \cdot 49$ on $5$
axes (cell-id-derived model\_family $\times$ task\_slice $\times$
$G \times$ temperature $\times$ seed is one design point), confirming
the manifests and the cells ledger describe the same factorial design
via parallel encodings (cell\_id string prefix vs structured
\texttt{cells.tsv} columns).

\textbf{Cross-paper coupling.} (i)~P5 iter-97 row 114 reported the
manifest-vs-cells \emph{schema} mismatch (3 of 5 actually-varying
axes absent as structured manifest keys); iter-105 closes the value
side: the $2$ remaining axes (model\_family, seed) are recoverable
from cell\_id parsing at $100\%$, so the recovered axis-set equals
the declared axis-set after a single parse. (ii)~P5 iter-93 row 109
reported bootstrap CIs on per-axis $\eta^{2}$ at $n{=}98$;
iter-105 confirms that the corpus design (factorial $98$) is exactly
the corpus on which iter-93's CIs are computed, so there is no
denominator inflation. (iii)~P6 iter-90 row 107 reported registry
cross-reference integrity on the zvf130 risk-index; iter-105's
audit-primitive finding (3/8 keys are constant) suggests the GRPO
Registry's audit should likewise split discriminative vs audit-primitive
keys.

\textbf{Operational recommendation.} Adopt the iter-105 split:
stack-discriminative fields ($5/8$) carry the stack axis; audit-primitive
fields ($3/8$) carry the run-time assurance. Report both
``$n_{\text{fields}}$ discriminative'' and ``$n_{\text{fields}}$
primitive'' on the MIN-REPORT coverage line, not the present-vs-absent
binary.

\textbf{Reproducibility.} Every artefact in
\texttt{platform_hybrid/experiments/results/p5p8/p5\_iter105\_*} is regenerated by the
single script under a fixed seed; the per-cell classification is exact
(JSON parse), and the bootstrap CIs are over a deterministic $n{=}98$
resampling (no randomness in the denominator).
